\documentclass[11pt]{article}

\usepackage[margin=1in]{geometry}
\usepackage{amsmath,amssymb,amsfonts}
\usepackage{bm}
\usepackage{hyperref}
\usepackage[numbers]{natbib}

\title{Multiplicity--Zeta--Phi--$\Pi$ Field Theory of Cognition:\\
Spectral $\varphi$--Invariants and a Preregistered EEG Test}
\author{Citizen Gardens \\ \\ The Foundation of Multiplicity}
\date{April 2, 2026}

\begin{document}
\maketitle

\begin{abstract}
Multiplicity--Zeta--Phi--$\Pi$ (MZ$\varphi\Pi$) Field Theory models human cognition as a multiplicative spectral dynamical system over a prime--indexed Hilbert space, with learning driven by resonance between internal and environmental eigenmodes.[file:1]
A central derived prediction is that self--organized cognitive dynamics exhibit golden--ratio ($\varphi \approx 1.618$) organization in their spectral peaks, especially in resting and low--pressure flow states.[file:1]
This report (i) summarizes the core mathematical architecture of MZ$\varphi\Pi$, (ii) refines the role of the golden ratio from universal invariant to dominant attractor under recursive growth, and (iii) presents a preregistered EEG study designed to test the $\varphi$--spectral prediction empirically.[file:1][web:18]
\end{abstract}

\section{Executive Summary}

MZ$\varphi\Pi$ Field Theory posits that cognitive states live in a prime--indexed Hilbert space $H = \ell^{2}(P)$, evolve under a multiplicative generator $\Pi$, and couple to the environment through a Dirichlet--type zeta transform that encodes resonance.[file:1]
Under a Fibonacci--like recursion on mode weights, the theory predicts that normalized growth ratios converge toward the golden ratio $\varphi$, making it a dominant scaling invariant of self--organizing cognitive fields.[file:1]
This scaling prediction, combined with prior work suggesting golden--ratio organization of EEG bands, motivates a concrete, preregistered test: adjacent dominant spectral peak ratios in resting EEG should approximate $\varphi$ and degrade under high pressure.[file:1][web:16][web:18]

\section{Mathematical Overview of MZ$\varphi\Pi$}

\subsection{Prime--Indexed Cognitive Space}

Let $P$ denote the set of prime numbers.
Cognitive states are modeled as elements of
\[
H = \ell^{2}(P) = \left\{ \bm{\lambda} = (\lambda_{p})_{p \in P} : \sum_{p \in P} |\lambda_{p}|^{2} < \infty \right\},
\]
with an orthonormal basis $\{\psi_{p}\}_{p \in P}$.[file:1]
At time $t$, the internal cognitive state and the environmental state are
\[
M(t) = \sum_{p \in P} \lambda_{p}(t)\,\psi_{p},\quad
E(t) = \sum_{p \in P} \mu_{p}(t)\,\phi_{p},
\]
where $\psi_{p},\phi_{p}$ are prime--indexed eigenmodes (e.g.\ phonological, visuospatial, narrative, social).[file:1]
The core Multiplicity axiom is that primes are the irreducible generators of recursive multiplicative structure, hence the natural index set for modes in systems that grow by multiplicative recursion.[file:1]

\subsection{Dirichlet--Spectral Transform and Resonance}

For a fixed (possibly finite) set of primes $P_{\max} \subset P$, define the \emph{cognitive Dirichlet transform}
\[
\zeta_{M}(s)
  = \sum_{p \leq P_{\max}} \lambda_{p}\,p^{-s},\qquad
\zeta_{E}(s)
  = \sum_{p \leq P_{\max}} \mu_{p}\,p^{-s},
\]
with $s = \sigma + i t$.[file:1]
In the full theory, $\zeta_{M}$ is treated as a formal Dirichlet series; convergence and analytic properties are asserted only for finite truncations used in empirical work.[file:1]
The \emph{resonance kernel} is
\[
R(s) = \zeta_{M}(s)\,\zeta_{E}(s)
      = \sum_{p,q} \lambda_{p}\mu_{q}\,(pq)^{-s},
\]
a multiplicative Dirichlet kernel encoding mode--mode coupling.[file:1]
Low primes dominate coarse--grained resonance, while higher primes encode finer sequential and symbolic structure.[file:1]

An effective learning rate can be written schematically as
\[
\frac{d M}{dt} = \alpha\,R_{\mathrm{eff}}(t),
\]
where $R_{\mathrm{eff}}(t)$ is a time--dependent functional derived from $R(s)$ in a task--relevant spectral window.[file:1]

\subsection{Recursive Growth and the Golden Ratio}

For each fixed prime mode $p$, MZ$\varphi\Pi$ assumes a recursive update under successful resonance episodes:
\begin{equation}
\lambda_{p}(t+1)
  = \lambda_{p}(t) + \lambda_{p}(t-1) + \epsilon_{p}(t),
\label{eq:fib}
\end{equation}
with bounded perturbations $\epsilon_{p}(t)$ capturing noise and minor nonlinearity.[file:1]
Under mild conditions on $\epsilon_{p}(t)$, the ratio
\[
\frac{\lambda_{p}(t+1)}{\lambda_{p}(t)} \to \varphi
  = \frac{1 + \sqrt{5}}{2}
\]
in expectation as $t \to \infty$, so $\varphi$ emerges as a dominant \emph{scaling} invariant for mode amplitudes.[file:1]

Substituting (\ref{eq:fib}) into the truncated Dirichlet transform yields a multiplicative cascade whose asymptotic scaling is controlled by $\varphi$.
This motivates the definition of a normalized spectral ratio functional
\[
M_{\varphi} = \lim_{s \to s_{0}}
  \frac{\zeta_{M}(s+1)}{\zeta_{M}(s)},
\]
which converges to $\varphi$ under Fibonacci--type growth.[file:1]
In this weakened form, $\varphi$ is treated as a \emph{dominant attractor} for a wide class of natural recursive trajectories, not as an exclusive universal invariant; other attractors may exist, but $\varphi$ is expected to be common.[file:1]

Independent empirical work supports golden--ratio organization of EEG frequency bands.
Klimesch and colleagues showed that canonical bands can be modeled as a geometric series with ratio close to $\varphi$.\cite{Pletzer2010,Klimesch2013}[web:16]
A recent large--sample study introduced the Phi Coupling Index (PCI) and found that human resting EEG theta--alpha ratios are significantly closer to $\varphi$ than to simple harmonics in $N=320$ participants.\cite{GoldenEEG2026}[web:18]

\subsection{Multiplicative Generator $\Pi$}

Cognitive evolution in $H$ is driven by a multiplicative operator
\[
(\Pi M)_{p} = \sum_{q} K_{pq}\,\lambda_{q},
\]
with coupling kernel
\[
K_{pq} = \eta(p,q)\,(pq)^{-s_{0}},
\]
where $\eta(p,q)$ is a multiplicity index counting recursive interactions of modes $p$ and $q$, and $s_{0}$ is a task--dependent spectral parameter.[file:1]
Theoretical work suggests an Euler--factorized form
\[
\Pi_{\mathrm{full}} = \bigotimes_{p \in P} \Pi_{p},\qquad
\Pi_{p} = \frac{1}{1 - \varphi\,p^{-s_{0}}},
\]
making $\Pi$ multiplicative over primes in analogy with the Euler product of the Riemann zeta function.[file:1]
For empirical purposes, a finite--dimensional \emph{$\Pi$--kernel} is estimated by truncating to low primes; this kernel is viewed as the observable shadow of the full $\Pi_{\mathrm{full}}$.[file:1]

Under repeated resonance--driven updates, the leading eigenvalue of empirically fitted $\Pi$--kernels is predicted to cluster near $|\lambda_{\max}|\approx \varphi$ across individuals, while eigenvectors differ by cognitive profile (e.g.\ dyslexic vs neurotypical).[file:1]

\section{Preregistered EEG Test of the $\varphi$--Spectral Prediction}

\subsection{Core Prediction}

MZ$\varphi\Pi$ predicts that in stable, high--resonance regimes such as eyes--closed rest and low--pressure flow states, adjacent dominant spectral peaks across canonical EEG bands approximate the golden ratio:
\[
\frac{f_{\alpha}}{f_{\theta}} \approx
\frac{f_{\beta}}{f_{\alpha}} \approx
\frac{f_{\gamma}}{f_{\beta}} \approx \varphi,
\]
with systematic degradation (increased deviation and variance) under high--pressure performance conditions.[file:1]
Empirically, golden--ratio organization has been observed in theta--alpha relationships and quantified by PCI, motivating a more general peak--ratio test.\cite{GoldenEEG2026}[web:18]

\subsection{Design Overview}

The preregistered study is a within--subject EEG experiment with four conditions: eyes--open rest, eyes--closed rest, low--pressure passive narrative, and high--pressure performance on matched narrative content.[file:1]
A between--subject factor distinguishes dyslexic and neurotypical participants, allowing tests of whether $\varphi$--organization is shared at rest across neurotypes.[file:1]

\paragraph{Participants.}
Target $N=36$ adults (18 with diagnosed dyslexia, 18 neurotypical controls), age 18--40, right--handed, no major neurological disorders.[file:1]
Dyslexia is confirmed by prior diagnosis and standardized reading measures; controls have no learning disability history and reading scores above a specified percentile.[file:1]

\paragraph{Conditions.}
\begin{itemize}
\item Eyes--open rest (5 min): fixation on a cross.
\item Eyes--closed rest (5 min).
\item Low--pressure passive: engaging narrative listening (10--15 min), explicitly no evaluation; post--condition flow ratings.
\item High--pressure performance: matched narrative with instructions emphasizing timed testing and evaluation; subsequent recall and pressure ratings.[file:1]
\end{itemize}

\paragraph{EEG Acquisition and Preprocessing.}
EEG is recorded with a 32--64 channel 10--20 system at $\ge 500$ Hz, bandpass 0.5--100 Hz, notch at 50/60 Hz.[file:1]
Preprocessing uses standard pipelines (MNE--Python/EEGLAB): re--referencing, bandpass 1--45 Hz, ICA for ocular/muscle artifacts, rejection of high--amplitude segments.[file:1]
Spectral analysis (Welch, 2 s windows, 50\% overlap) is computed over a posterior ROI (Pz, POz, Oz, O1, O2) where resting rhythms are robust.[file:1]

\paragraph{Peak Extraction.}
For each condition, the power spectral density (PSD) is smoothed (Gaussian kernel) and peaks are detected as local maxima with prominence $\ge 10\%$ of the condition maximum.[file:1]
Dominant peaks are identified per canonical band:
\[
\theta: 4\!-\!7\;\mathrm{Hz},\quad
\alpha: 8\!-\!12\;\mathrm{Hz},\quad
\beta: 13\!-\!30\;\mathrm{Hz},\quad
\gamma: 31\!-\!45\;\mathrm{Hz},
\]
and labeled $f_{\theta}, f_{\alpha}, f_{\beta}, f_{\gamma}$ when present.[file:1]

\subsection{Outcome Measures and Hypotheses}

Primary adjacent ratios are
\[
r_{\alpha\theta} = \frac{f_{\alpha}}{f_{\theta}},\quad
r_{\beta\alpha} = \frac{f_{\beta}}{f_{\alpha}},\quad
r_{\gamma\beta} = \frac{f_{\gamma}}{f_{\beta}},
\]
aggregated across participants and conditions.[file:1]
Secondary analyses use all detected peaks to compute data--driven adjacent ratios $r_{k} = f_{k+1}/f_{k}$ and fit $\varphi$--lattices $f_{k} \approx f_{1}\varphi^{k-1}$ via log--linear regression.[file:1]

The preregistered hypotheses are:

\begin{itemize}
\item H1 (rest $\varphi$--ratio): in eyes--closed rest, mean adjacent ratios lie within $[1.58,1.66]$; an equivalence test within this interval supports H1.[file:1]
\item H2 (pressure disruption): deviation $|r - \varphi|$ is significantly larger in high--pressure than in resting and low--pressure conditions.[file:1]
\item H3 (low--pressure flow): $\varphi$--clustering in low--pressure passive narratives is similar to or stronger than in rest, and correlates positively with self--reported flow.[file:1]
\item H4 (group similarity): dyslexic and neurotypical participants show equivalent $\varphi$--clustering at rest, consistent with a shared scaling attractor despite different task performance under pressure.[file:1]
\end{itemize}

\subsection{Falsification Criteria}

The preregistration specifies that the strong $\varphi$--specific claim is falsified if, in eyes--closed rest, the 95\% confidence or credible interval for the mean adjacent ratio excludes $[1.58,1.66]$ and no alternative narrow ratio band emerges consistently across individuals.[file:1]
Failure of the pressure manipulation to alter $\varphi$--clustering weakens the resonance--stress sub--claim even if golden--ratio organization is observed at rest.[file:1]

\section{Illustrative Analysis Code Snippet}

The following Python--style pseudocode illustrates how the spectral ratios will be computed from preprocessed EEG:

\begin{verbatim}
import mne
import numpy as np
from scipy.signal import welch, find_peaks

# epochs: MNE Epochs object for one condition and participant
roi_chs = ['Pz', 'POz', 'Oz', 'O1', 'O2']
data = epochs.copy().pick_channels(roi_chs).get_data()  # shape: (n_epochs, n_chs, n_times)

# Average over epochs and channels
avg = data.mean(axis=(0,1))

fs = epochs.info['sfreq']
freqs, psd = welch(avg, fs=fs, nperseg=int(2*fs), noverlap=int(1*fs))

# Smooth PSD (e.g., Gaussian convolution) omitted for brevity

# Peak detection
peaks, props = find_peaks(psd, height=0.1*psd.max())

# For each canonical band, pick the highest peak
bands = {'theta': (4,7), 'alpha': (8,12), 'beta': (13,30), 'gamma': (31,45)}
band_peaks = {}
for band, (fmin,fmax) in bands.items():
    idx = np.where((freqs[peaks] >= fmin) & (freqs[peaks] <= fmax))[0]
    if len(idx) > 0:
        # highest PSD peak in band
        best = idx[np.argmax(psd[peaks][idx])]
        band_peaks[band] = freqs[peaks][best]

# Compute ratios when possible
if all(b in band_peaks for b in ['theta','alpha']):
    r_alpha_theta = band_peaks['alpha'] / band_peaks['theta']
\end{verbatim}

This code is a condensed illustration; the preregistered pipeline includes full artifact handling, parameter choices, and statistical analysis as specified above.[file:1]

\section{Discussion and Outlook}

If the preregistered study finds robust $\varphi$--ratio organization in resting EEG and systematic disruption under pressure, it will provide strong evidence that MZ$\varphi\Pi$ has successfully predicted a nontrivial structural invariant of cognitive dynamics, beyond descriptive fitting.[file:1][web:18]
If the golden--ratio claim fails in this sharp form, the broader multiplicative framework may still hold but will require revision of its spectral invariants and attractor structure.[file:1]
In either case, the preregistration anchors the theory to falsifiable empirical commitments and informs subsequent work on $\Pi$--kernel estimation and pressure--resonance behavior.[file:1]

\bibliographystyle{plainnat}
\begin{thebibliography}{9}

\bibitem{Pletzer2010}
Pletzer, B., Kerschbaum, H., \& Klimesch, W.
\newblock When frequencies never synchronize: the golden mean and resting EEG.
\newblock \emph{Brain Research}, 1335, 2010.[web:16]

\bibitem{Klimesch2013}
Klimesch, W.
\newblock An algorithm for the frequency architecture of brain oscillations.
\newblock \emph{Brain and Cognition}, 2013.[web:16]

\bibitem{GoldenEEG2026}
[Authors].
\newblock Golden ratio organization in human EEG is associated with theta--alpha
frequency convergence.
\newblock \emph{Frontiers in Human Neuroscience}, 2026.[web:18]

\bibitem{MultiplicityMZPhiPi}
[Author].
\newblock Multiplicity--Zeta--Phi--$\Pi$ Field Theory.
\newblock Internal draft, 2026.[file:1]

\end{thebibliography}

\end{document}